\chapter{Dense Kernel Optimization and Zero-Heap RL}
\section{The Eradication of FxHashMap}
During the v1.2.0 optimization pass, DHAT (Dynamic Heap Analysis Tool) revealed that the primary source of remaining latency was the \texttt{FxHashMap} used for state lookups. While \texttt{FxHashMap} is fast, it relies on bucket probing. Hash collisions trigger unpredictable probing sequences, and load-factor thresholds trigger massive re-allocation spikes.

To achieve nanosecond stability, DTEAM eradicated \texttt{FxHashMap} from the hot path entirely, introducing the \textbf{Packed Key Table (PKT)} and \textbf{Dense Indexing}.

\section{The Packed Key Table}
\begin{definition}[Packed Key Table]
A contiguous vector of tuples $(H, K, V)$, where $H$ is a 64-bit FNV-1a hash. The vector is strictly sorted by $H$, and all lookups are performed via binary search.
\end{definition}

While binary search is theoretically $O(\log N)$ compared to the $O(1)$ average of a hash map, the PKT's contiguous memory layout guarantees perfect cache locality. More importantly, it \textit{never} allocates on insertion if pre-sized, and it \textit{never} suffers from stochastic probing jitter.

\section{Quantized State Calculus and the Copy Property}
In DTEAM v1.2.1, the RL state is defined entirely by 136 bits of stack-allocated data:
\begin{itemize}
    \item \textbf{Health/Metrics:} 8-bit integers ($i8$).
    \item \textbf{Marking Mask:} 64-bit bitset ($u64$).
    \item \textbf{Activity History:} 64-bit rolling FNV-1a hash ($u64$).
\end{itemize}
Because the state struct implements the Rust \texttt{Copy} trait, there are zero pointers, zero heap allocations, and zero cloning overheads when moving states between the Q-table and the agent policy.

\section{Empirical Validation: The 14.87ns Update}
By marrying the PKT with Quantized Copy States and returning borrowed slices rather than cloned vectors, DTEAM achieved a staggering performance milestone. As benchmarked via `divan`, a single Q-Learning update requires exactly \textbf{14.87 nanoseconds}. 

Furthermore, DHAT profiling confirmed that running 1,000,000 RL updates produces \textbf{zero steady-state heap allocations}. The RL update has been reduced to the level of a systems primitive.
